\documentclass[../Main.tex]{subfiles}
\begin{document}

\section{Discussion}
\label{chap:discussion}

In this project, we have had the opportunity to explore the genre of the alternative game. We have studied existing alternative games and invented new alternative game concepts. This section will discuss some of the choices and challenges that we encountered during the work.

The initial phase of the project focused on reviewing alternative games literature to get an introduction to the field. Reading alternative game articles also gave us inspiration for new alternative game concepts. This first phase of the project was essential to gain insight and give an introduction to the field of alternative games before we defined the project and context. We early on decided that we wanted to do something socially beneficial, and consequently, the concept we later would choose needed to reflect this decision.

The next step in the project was to perform a literature review of essential game development topics. This study was a central part of the preparations towards being able to develop a game, as both of us had minimal prior game development experience. Game design principles for designing an enjoyable, motivating, and engaging game were reviewed. This included studying GameFlow mechanisms, player types, and reward mechanisms to understand how to design enjoyable games for different players. 

During the literature review of alternative games, our goal was to understand the contributions of previous research. However, the literature review shed light on the lack of a concise definition of what an alternative game is. Thus we need a more precise definition and decided to define it ourselves. Consequently, the scope of the project was expanded to include defining what an alternative game is, which was not initially planned for. The result of our work became a definition of what an alternative game is and a framework that can be used to categorize existing concepts. Additionally, the model serves as a structured way to invent new concepts by combining different characteristics. This structured approach was useful later when we looked at existing alternative games and in the process of inventing new ones. The model and definition were based on our previous introduction and research of alternative games in addition to feedback and suggestions from our supervisor Alf Inge Wang. Because our definition is mainly based on our knowledge and experience in the field of alternative games, it would be naive to assume that it is complete. There may very well be additional aspects and considerations which we have not recognized.

After defining alternative games, the definition and model were used to study six existing alternative games to deepen our understanding of what alternative games are. The selection of games was based on their differences, where the purpose was to choose six distinct concepts. As mentioned above, the alternative games model was used to categorize the games to highlight their differences. The study revealed the diversity of technology in alternative games, such as input and output mechanisms.

During the period of researching existing alternative games, we realized how big the field is. The possibilities of alternative games are near endless. Consequently, we decided to limit our scope of potential ideas before trying to invent new concepts later in the project. Our decision became to focus on mobile applications based on the motivation of creating a socially beneficial concept and target as many users as possible. Smartphones have become extremely popular in the later years \citep{odea_cell_2020}. For this reason, a mobile application was a good option.

The choice to only consider mobile applications was helpful when inventing new concepts. The decision also shaped the content of the technology review, where we evaluated relevant technologies. During this process, we researched and explored central tools and frameworks. The knowledge obtained was useful later for choosing how to develop our chosen concept. The technology review also included a brief study of general technologies used in alternative games. It was relevant for us to get an overview of the field and not limit our scope too much in the beginning.

We consider the decision to focus on only mobile development as both positive and negative. In the early brainstorming session, we struggled with coming up with new concepts. The limitation made it a lot easier for us to invent new concepts by giving our minds a specific area to focus on, simplifying the brainstorming process. On the other hand, the decision meant giving up all potential alternative game concepts which were not mobile applications. However, we feel confident that our decision was “correct” based on our goal to target as many people as possible and our struggle with inventing concepts in the early stages.

After finishing the prestudy, our focus shifted to the task of coming up with, deciding, and concretizing a game concept. The first step was to continue the brainstorming process, as mentioned above. Our goal was to explore and invent multiple alternative game concepts, one of which we would continue working on and developing during our master’s thesis. We worked together with supervisor Wang to conceptualize ideas, and the result of the process became six different alternative game concepts. Supervisor Wang contributed with both suggestions for new concepts and improvement to the concepts we invented ourselves.

The concept which we decided to pursue is \textit{Hover}, a game where the player is awarded points for socially beneficial tasks. All the concepts were compared and considered against each other. Hover was chosen because it was the concept that was most in line with our motivation and the other considerations related to the development of our master’s thesis. We looked at how easy the concepts were to test, the number of potential test subjects, and how extendable the idea is. 

Due to the limited amount of time in the master’s thesis project, the chosen concept should be possible to translate into a Minimum Viable Product (MVP), while still including the possibility to improve the game experience and add more game elements. Hover was the game idea that overfall felt best suited. As a secondary consideration, Hover was also the concept that we were most excited about. We are eager to develop it and get feedback from real players to determine whether or not the game idea is exciting.

In the last part of the project, the focus was directed to further describing and deciding how we want to develop our chosen concept, Hover. The idea was polished, and we deliberated on game design and made suggestions for development tools and frameworks. The selection process here was, similarly to previous choices, based on the foundation and research we had previously done. 

To get a head start on the development process, before the master’s thesis part of the project, we spent some days putting together a quick prototype (illustrated and described in Chapter \ref{chap:chosen_tech}). The purpose was to test out the central components of the technological tools and frameworks we had chosen. Our initial impressions from this work are positive, and we are excited to continue working on this.

Even though issues and challenges certainly will appear, we feel confident that this specialization project has given us a solid foundation for our master’s thesis work. The next step now (further explained in Chapter \ref{chap:further_work}) is to plan how we want to develop the application and implement game design features. We are satisfied with our work so far and are eager to continue at the same pace to create the best possible end product in our master’s thesis.


\biblio % Needed for referencing to working when compiling individual subfiles - Do not remove
\end{document}